% Part: lambda-calculus
% Chapter: syntax
% Section: abbreviated-syntax

\documentclass[../../../include/open-logic-section]{subfiles}

\begin{document}

\olfileid[id]{lam}{syn}{abb}
\olsection{Sintaksis Singkat}

Suku-suku sebagaimana didefinisikan dalam \olref[trm]{defn:term} terkadang
merepotkan untuk ditulis, sehingga berguna untuk memperkenalkan sintaksis yang
lebih ringkas. Tentu saja, kita harus berhati-hati memastikan bahwa suku-suku
dalam notasi ringkas tersebut juga dapat dibaca secara unik. Salah satu versi
yang digunakan secara luas, yang disebut \emph{suku singkat}, adalah sebagai
berikut.

\begin{enumerate}
\item Apabila tanda kurung dihilangkan, aplikasi berlangsung dari kiri ke
  kanan. Sebagai contoh, jika $M$, $N$, $P$, dan~$Q$ adalah suku, maka $MNPQ$
  menyingkat $(((MN)P)Q)$.
\item Sekali lagi, apabila tanda kurung dihilangkan, abstraksi lambda diberi
  lingkup seluas mungkin. Sebagai contoh, $\lambd[x][MNP]$ dibaca sebagai
  $(\lambd[x][((MN)P)])$.
\item Satu lambda dapat digunakan untuk mengabstraksikan beberapa variabel.
  Sebagai contoh, $\lambd[xyz][M]$ menyingkat
  $\lambd[x][\lambd[y][\lambd[z][M]]]$.
\end{enumerate}

Sebagai contoh,
\[
\lambd[xy][xxyx \lambd[z][xz]]
\]
menyingkat
\[
(\lambd[x][(\lambd[y][((((xx)y)x)(\lambd[z][(xz)]))])]).
\]

\begin{prob}
Uraikan suku singkat $\lambd[g][(\lambd[x][g (x x)])
  \lambd[x][g (x x)]]$ ke bentuk lengkapnya.
\end{prob}

\end{document}
